\section{Recommendations and Conclusion}
\label{header:recs}

\paragraph{}The recommendations developed in this section follow directly from the empirical and interpretive results of Sections~\ref{header:results} and~\ref{header:discussion}. They are organised from the most operationally immediate (adaptive QR sizing, in response to the A/B finding) through the most architecturally consequential (post-quantum migration), and each is grounded either in a participant-elicited concern or in a limitation flagged in Section~\ref{sect:disc_limits}. The section concludes with a brief synthesis tying the prototype's results back to the gaps identified in the related literature.

\subsection{Adaptive QR Sizing by Transaction Risk}
\label{sect:rec_adaptive}

\paragraph{}The most direct response to the A/B finding of Section~\ref{sect:results_uat_qr} is to make the QR variant a per-transaction policy decision keyed by transaction value, rather than a system-wide configuration. The protocol already accepts transaction-value-keyed policy in the threshold itself --- $t = 1$ for low-value transactions, $t \geq 2$ for high-value transactions --- and the same Policy Engine (Section~\ref{sect:meth2_b}, item~d) can be extended to also choose the QR encoding strategy:

\begin{itemize}
  \item \textbf{Low-value transactions.} Use the Small QR variant (reference token only). The transactions on this path are below the threshold for cryptographic multi-party authorisation by design; the marginal security cost of routing the verification envelope through the proxy's retrieval channel is therefore commensurate with the value at risk, and the user receives the instantaneous scan that the UAT validated for Group~B.
  \item \textbf{High-value transactions.} Use the Big QR variant (full in-band payload). Here the transaction value justifies the additional self-containment property; the user has already accepted multi-second latency for approver coordination, so a 2--3 second autofocus penalty is a small additional fraction of the total transaction time, and not a marginal cost on top of an otherwise instant flow.
\end{itemize}

\paragraph{}A complementary refinement, drawing on the literature on QR autofocus latency~\cite{corpuz2025} and on the structural properties of Reed--Solomon error correction~\cite{reed1960, bajaj2024}, is to lower the error correction level for the Big QR variant from H (30\%) to L (7\%) or M (15\%) in indoor, well-lit settings, freeing module budget to encode the same payload in a lower QR version with a sparser grid. The error correction level can itself be elevated dynamically by the Policy Engine in conditions where scan reliability is degraded (e.g., outdoor or low-light merchant terminals), provided such conditions can be inferred from the verifier's scanner telemetry.

\subsection{Joint Fund Replenishment (``Deposit'') Feature}
\label{sect:rec_deposit}

\paragraph{}Several UAT participants observed that the prototype currently models only the spending side of a joint account: the protocol authorises debits via threshold approval, but there is no symmetric flow for credits (Section~\ref{sect:results_uat_open}). Adding a deposit feature is straightforward at the implementation level --- the \texttt{GroupBalance} document already supports increments as well as decrements, and the existing audit log schema accommodates credit events --- but the design question is whether deposits should themselves require multi-party authorisation, and if so, on what basis.

\paragraph{}Three policy designs are worth distinguishing. The first, \textit{permissive deposits}, treats credits as benign by default: any member may deposit funds into the joint balance without quorum approval, on the reasoning that an addition cannot harm other members. The second, \textit{symmetric thresholding}, applies the same value-keyed threshold to deposits as to withdrawals, on the reasoning that very large deposits (e.g., a member retitling external funds into the joint account, or a merchant refund of a previously-authorised purchase) should be visible to and acknowledged by the same quorum. The third, \textit{asymmetric thresholding}, applies a strictly lower threshold to deposits than to withdrawals (for instance, $t_{\text{deposit}} = 1$ regardless of value), on the reasoning that the trust placed in deposits is structurally different from the trust placed in withdrawals.

\paragraph{}The recommendation of the present study is the second policy --- symmetric thresholding --- with two additional safeguards. First, the deposit message format should include an explicit \textit{source-of-funds} field, distinguishing user deposits (from a member's own external wallet) from inbound transfers (from a merchant refund or from a third party); the threshold should depend on this field as well as on the amount. Second, all deposits, regardless of threshold, should generate a non-suppressible audit event visible to every member, so that even a $t = 1$ permissive deposit is observable by the rest of the quorum.

\subsection{Multi-Factor Approval Confirmation}
\label{sect:rec_mfa}

\paragraph{}The second recurring suggestion in the UAT was that the protocol should require an explicit credential confirmation at the point of approval, rather than treating possession of an unlocked phone as constitutive of consent (Section~\ref{sect:results_uat_open}). The recommendation is to bind every \texttt{submit\_share} event to a fresh per-approval authentication artefact, drawn from one of three categories:

\begin{itemize}
  \item \textbf{Knowledge factor.} A short PIN (4--6 digits) or a passphrase entered at the point of approval. The PIN is held client-side and is never transmitted; instead, it is used as the input to a key-derivation function that unlocks the on-device storage holding the participant's BLS private share. A failed PIN entry leaves the share inaccessible.
  \item \textbf{Inherence factor.} A platform biometric (fingerprint or Face~ID) gating the same on-device unlock. The biometric never leaves the device; the platform's secure enclave returns a binary attestation that the biometric matched, and that attestation gates the unlock of the share. This is the strongest of the three categories with respect to phishing and shoulder-surfing, although it depends on the platform's biometric implementation and inherits its failure modes.
  \item \textbf{Possession factor.} A hardware token (e.g., a YubiKey or an equivalent FIDO2 authenticator) presented at the point of approval. The token holds the BLS share or a wrapping key for the share, and the per-approval signing operation cannot complete without the token's presence. This is the strongest defence against device-loss scenarios but introduces a hardware dependency.
\end{itemize}

\paragraph{}The recommendation of the present study is to support \textit{any} of the three at the participant's discretion --- mirroring the standard FIDO2 model --- with biometric gating as the default on platforms that provide a hardware-backed secure enclave. The protocol need not be aware of which factor was used, since the authentication is local to the participant's device; what changes is the semantics of a successful approval event, which now reads ``the participant intentionally authorised this transaction with a fresh credential check'' rather than ``a process running on the participant's phone produced a partial signature''. This is the kind of visible-procedural ratification of an invisible-cryptographic step that Section~\ref{sect:disc_perceived} identified as load-bearing for perceived security.

\subsection{Proactive Secret Sharing and Share Refresh}
\label{sect:rec_proactive}

\paragraph{}Section~\ref{sect:disc_limits} flagged the endpoint-compromise scenario --- a coalition of $t$ compromised devices producing a forged authorisation --- as out of scope for the present evaluation. The standard cryptographic defence against that scenario is \textit{proactive secret sharing}, in which the private shares are periodically re-randomised (without changing the underlying secret or the joint public key) so that shares captured before a refresh become useless after it~\cite{shamir1979, gennaro2007}. In a payment-system context, a natural refresh schedule is monthly or quarterly, with an additional emergency refresh triggered by any anomaly in a participant's signing pattern (e.g., signatures from unfamiliar locations or unusual amounts).

\paragraph{}A practical refresh protocol in the present setting would proceed in four steps: (i) the participants jointly run a fresh distributed key-generation round on the same polynomial degree, producing a new sharing of the same secret; (ii) each participant's old share is overwritten with the new share inside its on-device protected storage; (iii) any in-flight transactions that were signed under the old shares complete normally, since the joint public key is unchanged; and (iv) the proxy logs the refresh as an audit event but learns nothing about the new shares. The cost is a coordination round that is roughly equivalent to the cost of a single high-value transaction signing, which is small enough that monthly refreshes would not impose a meaningful operational burden.

\paragraph{}A complementary recommendation is that the threshold $t$ itself be re-derivable: under the present design, $t$ is fixed at account creation and parameterises the entire share polynomial. In production, it should be possible to issue a new share polynomial of a different degree (and therefore a different threshold) without forcing every member to discard their existing wallet identifier. The standard construction here is the resharing of a known public key under a new polynomial, which is well-established in the threshold-cryptography literature~\cite{gennaro2007, komlo2021}.

\subsection{Hardware-Backed Key Storage}
\label{sect:rec_hardware}

\paragraph{}The current prototype stores private shares in process memory on the server-side simulation. In a production deployment, the shares should reside in hardware-backed storage on each participant's device --- specifically, in a secure enclave or trusted execution environment, with the share never exposed to the application process even during signing. On modern mobile platforms, this is supported by the Secure Enclave (Apple) and the StrongBox-backed Keystore (Android); on desktop and server platforms, it is supported by hardware security modules and by attestable confidential-compute environments.

\paragraph{}Two consequences follow from this recommendation. First, the partial-signing operation must be exposed as an opaque service from within the secure enclave: the application can submit a digest $H(M)$ and receive a partial signature $\sigma_i$, but cannot read $k_i$ itself. Second, the secure enclave's attestation primitives can be exposed to the verifier, allowing the verifier to gate transaction acceptance on a proof that the partial signatures it received were produced inside genuine enclaves. This is a strictly stronger property than the present prototype provides, and addresses the endpoint-compromise scenario at the hardware level rather than the protocol level.

\subsection{Toward Post-Quantum Threshold Signatures}
\label{sect:rec_pqc}

\paragraph{}BLS12-381's security depends on the hardness of the elliptic-curve discrete-logarithm problem, which is broken in polynomial time by a sufficiently large quantum computer running Shor's algorithm. The current consensus among standardisation bodies is that this is a horizon risk rather than an immediate one, but for a payment system that may be deployed for decades the migration path should be considered now rather than at the moment a cryptographically relevant quantum computer is announced.

\paragraph{}Several post-quantum threshold-signature constructions are currently active in the research literature, including lattice-based threshold variants of Dilithium and hash-based threshold variants drawing on SPHINCS+. None has yet reached the production-readiness or the signature-compactness of BLS, and the research-grade post-quantum threshold signatures are typically an order of magnitude larger in signature size, which would interact unfavourably with the QR data-density observation of Section~\ref{sect:disc_qr}. The recommendation of the present study is therefore to design the protocol's signature-handling logic in a primitive-agnostic way: the verifier should treat $F_{\text{sig}}$ as an opaque binary blob whose verification predicate is supplied by an interchangeable primitive module, so that a future migration to a post-quantum primitive does not require touching the application logic, only the primitive module. This is in keeping with the broader architectural principle that cryptographic primitives should be interchangeable behind a stable interface, which the present prototype's separation between \texttt{utils/cryptoUtils.js} and \texttt{controllers/socketHandler.js} already approximates.

\subsection{Conclusion}
\label{sect:conclusion}

\paragraph{}This study set out to investigate whether a $(t, n)$ threshold-signature scheme --- specifically, BLS12-381 partial signatures aggregated under Lagrange interpolation and verified by a bilinear pairing --- could close three persistent gaps identified in the QR-payment security literature: the single-key failure mode of conventional dynamic-QR schemes, the post-hoc nature of machine-learning-based QR-payload defences, and the absence of multi-party authorisation primitives that integrate naturally with payment flows. Section~\ref{header:rrl} mapped those gaps; Section~\ref{header:methodology} proposed a hybrid construction in which threshold signatures, dynamic QR generation, ephemeral session keys, and a non-key-bearing proxy collectively address them; Section~\ref{header:results} evaluated the resulting prototype against four orthogonal acceptance criteria.

\paragraph{}The cryptographic claims hold up. Every legitimate 3-of-5 aggregate verifies; no 2-of-5 partial aggregate verifies; no payload mutation goes undetected; no off-curve point is admitted. The performance claims also hold up: every one of the 30 latency trials completed under 200~ms, with a 99th-percentile margin of more than $25\times$. The adversarial claims hold up: thirteen distinct attack scenarios, spanning offline cryptographic forgery and online protocol manipulation, were either correctly blocked or, in the two positive controls, correctly accepted. The user-acceptance claims hold up: ten participants in a between-subjects A/B study completed all ten task scenarios in both groups without failure, returned strongly positive scores on a six-item adaptation of the System Usability Scale, and unanimously rated the multi-party approval workflow as making the system safer than a single-user equivalent.

\paragraph{}The principal qualitative caveat that emerged from the evaluation is the QR-recognition asymmetry between the Big QR and Small QR variants, which shifts the design conversation from ``does the protocol work'' to ``how should the protocol's encoding be tuned to the transaction's risk profile''. This is not a refutation of the protocol's central claims; it is an engineering refinement, and it admits a clean policy-driven solution (Section~\ref{sect:rec_adaptive}) that exploits the same value-keyed Policy Engine that the protocol already uses to choose its threshold. The two participant-elicited feature requests --- a deposit / fund-replenishment flow, and stricter credential confirmation at the point of approval --- map onto two recommendations (Sections~\ref{sect:rec_deposit} and~\ref{sect:rec_mfa}) that are likewise additive to the existing protocol rather than transformative of it.

\paragraph{}Taken together, the cryptographic, performance, adversarial, and user-acceptance results jointly support the claim that threshold cryptography, when integrated carefully with dynamic QR generation and a non-key-bearing proxy relay, is a practical primitive for joint-account payment systems --- not merely a theoretical one. The remaining work, surveyed in Sections~\ref{sect:rec_proactive} through~\ref{sect:rec_pqc}, concerns the long-horizon questions of share refresh, hardware-backed storage, and post-quantum migration, none of which is unique to the present study but each of which is rendered tractable by the architectural separation between cryptographic primitive and application logic that the prototype establishes.
